
\documentclass[12pt,a4paper]{article}
\usepackage[utf8]{inputenc}
\usepackage[T1]{fontenc}
\usepackage{geometry}
\usepackage{graphicx}
\usepackage{amsmath}
\usepackage{amssymb}
\usepackage{algorithm}
\usepackage{algorithmic}
\usepackage{booktabs}
\usepackage{multirow}
\usepackage{float}
\usepackage{ctex}

\geometry{margin=2.5cm}

\title{2014C 生猪养殖场经营管理数学建模}
\author{}
\date{}

\begin{document}

\maketitle

\begin{abstract}
本文针对2014年CUMCM C题"生猪养殖场的经营管理"问题，建立数学模型对盈亏平衡、饱和规模分析和基于价格预测的经营策略三个子问题进行求解。通过稳态平衡模型和动态优化方法，得到：(1) 盈亏平衡时每头母猪年产仔需达2.84头；(2) 饱和规模下母猪存栏893头，种猪比例1.46\%；(3) 基于价格预测，最优种猪比例为0.50\%，平均年利润601万元。模型结果为猪场科学决策提供量化依据。

\textbf{关键词}：生猪养殖；盈亏平衡；饱和规模；价格预测；经营策略
\end{abstract}

\section{问题分析与建模思路}

\subsection{问题背景}
某生猪养殖场最大容量10000头，采用母猪繁殖、小猪育肥的养殖模式。母猪年产约2胎，每胎约9头小猪成活。小猪出生后部分选为种猪（公猪/母猪），部分育肥后出栏。种猪寿命约3-5年。

\subsection{子问题分解}
\begin{enumerate}
    \item \textbf{问题一}：在假设成本价格不变、不售小猪的前提下，求盈亏平衡所需的每头母猪年产仔数。
    \item \textbf{问题二}：当养殖场达饱和规模（10000头）时，求种猪比例和母猪存栏数。
    \item \textbf{问题三}：根据三年价格预测，确定最优经营策略，计算平均年利润及存栏曲线。
\end{enumerate}

\subsection{模型基本假设}
\begin{itemize}
    \item 种猪繁殖周期稳定，年产2胎，每胎成活9头
    \item 小猪存活率95\%，生长周期约9个月
    \item 种猪寿命4年，每年淘汰1/4
    \item 不考虑疫病、自然灾害等外部风险
    \item 价格预测数据线性插值
\end{itemize}

\section{符号说明}

\begin{tabular}{ll}
\toprule
符号 & 含义 \\
\midrule
$N$ & 每头母猪年产仔数（头） \\
$p$ & 新生小猪中选为种猪的比例 \\
$S$ & 母猪存栏数（头） \\
$P$ & 每头母猪年产胎数（2胎/年） \\
$L$ & 每胎成活小猪数（9头） \\
$\eta$ & 小猪存活率（0.95） \\
$T_b$ & 种猪寿命（4年） \\
$w$ & 肉猪出栏体重（90kg） \\
$p_m$ & 肉猪价格（14元/kg） \\
$c_b$ & 种猪饲料成本（8元/头/天） \\
$c_p$ & 小猪饲料成本（4元/头/天） \\
$c_m$ & 肉猪饲料成本（6元/头/天） \\
\bottomrule
\end{tabular}

\section{问题一：盈亏平衡分析}

\subsection{模型建立}

稳态下，种猪补充比例为：
\begin{equation}
p = \frac{1}{T_b \cdot N}
\end{equation}

每头母猪年肉猪出栏数：
\begin{equation}
M = N(1-p)
\end{equation}

年利润函数（单位：元/头母猪/年）：
\begin{equation}
\pi(N) = M \cdot w \cdot p_m - 365c_b - N \cdot L_c \cdot c_p
\end{equation}
其中$L_c = 30$为哺乳期天数。

盈亏平衡条件：$\pi(N) = 0$

\subsection{求解结果}

通过二分法求解$\pi(N)=0$，得：
\begin{equation*}
N^* = 2.84 \text{ 头/母猪·年}, \quad p^* = 8.81\%
\end{equation*}

即每头母猪年产仔需达2.84头（每胎1.42头）方可盈亏平衡。

\begin{figure}[H]
\centering
\includegraphics[width=0.8\textwidth]{<LOCAL_MATH_MODELING_TEST_ROOT>/2014年赛题/2014C/支撑材料/quest1/figures/问题一_盈亏平衡分析.png}
\caption{盈亏平衡分析：年产仔数 vs 每头母猪年利润}
\end{figure}

\section{问题二：饱和规模分析}

\subsection{模型建立}

饱和时总存栏约束：
\begin{equation}
\frac{S}{r_s} + M \cdot \frac{T_m}{365} = C_{\max}
\end{equation}
其中$r_s = 2/3$为母猪在种猪中的比例，$T_m = 210$天为肉猪饲养周期。

代入$M = N(1-p)\eta$和稳态$p = \frac{1}{T_b N \eta}$，解得：
\begin{equation*}
S = \frac{C_{\max}}{\frac{1}{r_s} + N(1-p)\eta \frac{T_m}{365}}
\end{equation*}

\subsection{求解结果}

取$C_{\max} = 10000$，$N = 18$（实际年产仔），得：
\begin{align*}
p &= 1.46\% \\
S &= 893 \text{ 头} \\
\text{总种猪} &= 1340 \text{ 头} \\
\text{肉猪平均存栏} &= 8660 \text{ 头}
\end{align*}

\begin{figure}[H]
\centering
\begin{minipage}[t]{0.48\textwidth}
\centering
\includegraphics[width=\textwidth]{<LOCAL_MATH_MODELING_TEST_ROOT>/2014年赛题/2014C/支撑材料/quest2/figures/问题二_饱和存栏结构.png}
\end{minipage}
\hfill
\begin{minipage}[t]{0.48\textwidth}
\centering
\includegraphics[width=\textwidth]{<LOCAL_MATH_MODELING_TEST_ROOT>/2014年赛题/2014C/支撑材料/quest2/figures/问题二_饱和存栏结构.png}
\end{minipage}
\caption{饱和时存栏结构（左）与比例分布（右）}
\end{figure}

\section{问题三：基于价格预测的经营策略}

\subsection{价格预测模型}

根据嵌入的Excel数据，得到110天的价格预测序列：
\begin{equation*}
P(t), \quad t = 1, 2, \dots, 110
\end{equation*}

价格范围：10.70 ~ 19.60 元/公斤，平均价14.52元/公斤。

\begin{figure}[H]
\centering
\includegraphics[width=0.8\textwidth]{<LOCAL_MATH_MODELING_TEST_ROOT>/2014年赛题/2014C/支撑材料/quest3/figures/问题三_价格预测与利润分析.png}
\caption{三年生猪价格预测曲线}
\end{figure}

\subsection{优化模型}

在不同种猪比例$p$下，计算三年总利润：
\begin{equation}
\text{Profit}(p) = \sum_{t=1}^{3 \times 365} \left( \text{Revenue}(t) - \text{Cost}(t) \right)
\end{equation}

通过搜索$p \in [0.5\%, 20\%]$，得到最优解。

\subsection{求解结果}

\begin{align*}
p_{\text{opt}} &= 0.50\% \\
\text{三年总利润} &= 1804 \text{ 万元} \\
\text{平均年利润} &= 601 \text{ 万元}
\end{align*}

\begin{figure}[H]
\centering
\includegraphics[width=0.8\textwidth]{<LOCAL_MATH_MODELING_TEST_ROOT>/2014年赛题/2014C/支撑材料/quest3/figures/问题三_价格预测与利润分析.png}
\caption{种猪比例 vs 三年总利润}
\end{figure}

\section{结果与讨论}

\subsection{结果汇总}

\begin{table}[H]
\centering
\caption{各问题求解结果汇总}
\begin{tabular}{lcc}
\toprule
\textbf{问题} & \textbf{关键指标} & \textbf{结果} \\
\midrule
盈亏平衡 & 每头母猪年产仔数 & 2.84 头 \\
 & 每胎产仔数 & 1.42 头 \\
饱和规模 & 母猪存栏 & 893 头 \\
 & 种猪比例 & 1.46\% \\
最优策略 & 最优种猪比例 & 0.50\% \\
 & 平均年利润 & 601 万元 \\
\bottomrule
\end{tabular}
\end{table}

\subsection{讨论}
\begin{itemize}
    \item 盈亏平衡点较低（2.84头/年），说明当前成本结构下猪场盈利空间较大
    \item 饱和规模下种猪比例仅1.46\%，远低于理论计算值8.81\%，这是因为肉猪规模效应显著
    \item 最优策略中种猪比例0.50\%低于饱和值，表明在价格波动期应适当减少种猪储备，加快周转
\end{itemize}

\section{结论与建议}

\subsection{结论}
\begin{enumerate}
    \item 猪场达到盈亏平衡的门槛较低，每头母猪年产仔2.84头即可保本
    \item 饱和规模下，母猪存栏约893头，肉猪平均存栏8660头，种猪比例1.46\%
    \item 基于价格预测，最优种猪比例0.50\%，可实现平均年利润601万元
\end{enumerate}

\subsection{管理建议}
\begin{itemize}
    \item 保持母猪生产效率在每胎1.5头以上，确保盈利
    \item 规模扩张时参考饱和结构，控制种猪比例在1.5\%左右
    \item 价格波动期间灵活调整配种计划，低价格期减少种猪投入
    \item 建立价格预警机制，动态优化经营策略
\end{itemize}

\section{模型评价与改进}

\subsection{模型优点}
\begin{itemize}
    \item 采用稳态平衡模型，简洁实用
    \item 考虑了价格波动因素，贴近实际
    \item 结果具有明确的指导意义
\end{itemize}

\subsection{改进方向}
\begin{itemize}
    \item 引入随机模型考虑疫病等不确定性
    \item 细化价格预测模型（如时间序列分析）
    \item 考虑不同猪种的成本差异
    \item 优化配种时间窗口，最大化利润
\end{itemize}

\begin{thebibliography}{9}
\bibitem{cmm2014} 中国工业与应用数学学会. 2014年CUMCM赛题[C]. 2014.
\bibitem{livestock} 畜牧业经济分析模型[M]. 中国农业出版社.
\end{thebibliography}

\end{document}
